% ===========================================================================
% CHAPTER 4
% ===========================================================================

\chapter{System Architecture and Data Design}
\label{chap:architecture-data-design}

\section{Architectural Overview and Data Flow}
\label{sec:architectural-overview}

SSL Guardian follows a layered, data-centric architecture that separates data
sources, certificate acquisition and maintenance, persistent data, analytical
processing, application services, and presentation. The persistent datasets
form the principal exchange boundary between these subsystems, allowing each
stage to perform a distinct architectural role. Figure~\ref{fig:overall-architecture}
presents the overall organization and high-level data flow.

\begin{figure}[H]
  \centering
  \begin{tikzpicture}[
    font=\small,
    >=Latex,
    source/.style={draw, dashed, rounded corners=3pt, align=center, text width=3.0cm, minimum height=0.92cm, line width=0.65pt},
    process/.style={draw, rounded corners=3pt, align=center, text width=3.45cm, minimum height=1.02cm, fill=black!5, line width=0.7pt},
    store/.style={draw, rounded corners=3pt, align=center, text width=3.35cm, minimum height=1.02cm, fill=black!12, line width=0.9pt},
    optional/.style={draw, dashed, rounded corners=3pt, align=center, text width=3.05cm, minimum height=0.92cm, fill=black!2, line width=0.65pt},
    flow/.style={-{Latex[length=2.4mm]}, line width=0.7pt},
    exchange/.style={{Latex[length=2.4mm]}-{Latex[length=2.4mm]}, line width=0.7pt}
  ]
    \node[source] (domains) at (-5.25,7.0) {Domain\\Datasets};
    \node[source] (tls) at (-1.7,7.0) {Internet TLS\\Endpoints};
    \node[source] (ct) at (3.8,7.0) {Certificate\\Transparency Logs};

    \node[process] (acquisition) at (-3.8,5.25) {Domain Certificate\\Acquisition and Processing};
    \node[process] (maintenance) at (3.8,5.25) {CT-Assisted Dataset\\Maintenance};

    \node[store] (primary-data) at (-3.8,3.25) {Primary Certificate\\Dataset};
    \node[process] (analytics) at (0,3.25) {Analytics\\Processing};
    \node[store] (results) at (3.8,3.25) {Analytical Results\\Dataset};

    \node[process] (services) at (0,1.35) {Application\\Services};
    % \node[optional] (cache) at (4.15,1.35) {Optional Response\\Cache};
    \node[process] (dashboard) at (0,-0.35) {Interactive\\Dashboard};
    \node[source] (users) at (4.15,-0.35) {Platform\\Users};

    \draw[dashed, rounded corners=4pt, line width=0.65pt]
      (-5.9,-1.75) rectangle (5.85,-3.55);
    \node[font=\small\bfseries] at (0,-1.48)
      {Independent Experimental Path};
    \node[source] (ip-input) at (-4.0,-2.65) {IP Address\\Ranges};
    \node[process] (ip-acquisition) at (0,-2.65) {Experimental IP\\Certificate Acquisition};
    \node[store] (ip-data) at (4.0,-2.65) {Separate IP\\Dataset};

    \draw[flow] (domains) -- (acquisition);
    \draw[flow] (tls) -- (acquisition);
    \draw[flow] (ct) -- (maintenance);
    \draw[flow] (maintenance.west) -- node[above, font=\scriptsize, align=center]
      {new and renewal\\domain candidates} (acquisition.east);
    \draw[flow] (acquisition) -- (primary-data);
    \draw[flow] (primary-data) -- (analytics);
    \draw[flow] (analytics) -- (results);
    \draw[flow] (primary-data) -- (services);
    \draw[flow] (results) -- (services);
    % \draw[exchange] (services) -- (cache);
    \draw[flow] (services) -- (dashboard);
    \draw[exchange] (dashboard) -- (users);
    \draw[flow] (ip-input) -- (ip-acquisition);
    \draw[flow] (ip-acquisition) -- (ip-data);
  \end{tikzpicture}
  \caption{End-to-end Architecture of the SSL/TLS Certificate Analytics Platform}
  \label{fig:overall-architecture}
\end{figure}

The acquisition layer follows two coordinated paths. Domain datasets provide
targets for certificate retrieval from Internet TLS endpoints, while
Certificate Transparency observations identify new and renewal candidates for
the same acquisition capability. Both paths therefore converge on the primary
certificate dataset. The experimental IP-based acquisition path remains
architecturally separate, and its dataset is not consumed by the primary
analytics or dashboard workflow.

The analytical layer transforms primary certificate data into global and
country-level analytical results. The service layer accesses both primary
records and analytical results, with an optional cache alongside it for selected
responses. The presentation layer obtains this information through the
application services and exposes it through the interactive dashboard, which
forms the system boundary presented to users.

\section{Certificate Data and Scope Model}
\label{sec:mongodb-data-model}

Each primary certificate record contains the parsed X.509 object together with
attributes introduced during acquisition. Table~\ref{tab:main-certificate-fields}
summarizes the principal field groups used by the analytical platform.

\begin{table}[H]
  \centering
  \small
  \setlength{\tabcolsep}{4pt}
  \caption{Principal certificate data fields}
  \label{tab:main-certificate-fields}
  \begin{tabularx}{\textwidth}{@{}>{\RaggedRight\arraybackslash}p{0.23\textwidth}>{\RaggedRight\arraybackslash}p{0.29\textwidth}Y@{}}
    \toprule
    \textbf{Field group} & \textbf{Representative content} & \textbf{Analytical purpose} \\
    \midrule
    Acquisition context & Domain, derived scope, scan time, and leaf classification & Search, scope filtering, and acquisition tracking \\
    Certificate identity & Serial number and certificate fingerprints & Lookup, deduplication, and certificate correlation \\
    Validity & Start time, end time, and lifetime & Current status, expiry windows, and validity distributions \\
    Issuer and subject & Distinguished names and organization attributes & Identity, issuer, and CA analysis \\
    Public key & Algorithm, key length, and SPKI fingerprint & Cryptographic distributions and shared-key detection \\
    Signature & Signature algorithm and hash information & Algorithm, hash-strength, and compliance analysis \\
    Extensions & SANs, key usages, policies, constraints, AIA, and CRL data & SAN analysis, usage analysis, and CA policy scoring \\
    Certificate-quality findings & ZLint flags and issue categories & Quality grading, risk indicators, and CA comparison \\
    \bottomrule
  \end{tabularx}
\end{table}

The platform maintains three related logical data groups, illustrated in
Figure~\ref{fig:logical-data-model}. Processing and staging data retains CT
observations, acquisition state, temporary results, and recovery metadata; it
supports CT-assisted maintenance but is not a mandatory precursor to direct
domain acquisition. Primary certificate records provide the detailed
analytical source, while analytical results store scope-tagged summaries and
shared-key groups derived from that source.

\begin{figure}[H]
  \centering
  \begin{tikzpicture}[
    font=\scriptsize,
    >=Latex,
    source/.style={draw, dashed, rounded corners=2pt, align=center, text width=1.65cm, minimum height=0.65cm, inner sep=3pt, line width=0.55pt},
    process/.style={draw, rounded corners=2pt, align=center, text width=1.65cm, minimum height=0.82cm, inner sep=3pt, fill=black!3, line width=0.6pt},
    data/.style={draw, rounded corners=2pt, align=center, text width=3.1cm, minimum height=2.45cm, fill=black!7, inner sep=3pt, line width=0.65pt},
    primary/.style={data, very thick, fill=black!11},
    flow/.style={-{Latex[length=2.0mm]}, line width=0.6pt}
  ]
    \node[data] (staging) at (-5.55,0) {\textbf{Processing \& Staging Data}\\[2pt]
      CT observations/candidates\\
      Task state\\
      Temporary results\\
      Recovery/run metadata};

    \node[process] (acquisition) at (-2.75,0) {Certificate\\Acquisition};

    \node[primary] (certificates) at (0,0) {\textbf{Primary Certificate Records}\\[2pt]
      Acquisition context and scope\\
      Parsed X.509 data\\
      Certificate/key fingerprints\\
      Certificate-quality findings};

    \node[process] (analytics) at (2.75,0) {Analytics\\Processing};

    \node[data] (results) at (5.55,0) {\textbf{Analytical Results}\\[2pt]
      Scoped summaries\\
      Certificate analytics\\
      CA/shared-key analytics\\
      Computation metadata};

    \node[source] (domain-inputs) at (-2.75,1.75) {Domain Inputs};

    \draw[flow] (staging.east) -- (acquisition.west);
    \draw[flow] (domain-inputs.south) -- (acquisition.north);
    \draw[flow] (acquisition.east) -- (certificates.west);
    \draw[flow] (certificates.east) -- (analytics.west);
    \draw[flow] (analytics.east) -- (results.west);
  \end{tikzpicture}
  \caption{Logical Organization of Certificate and Analytical Data}
  \label{fig:logical-data-model}
\end{figure}

The \texttt{scope} field stores the final label of the acquired domain; it can
therefore contain a generic TLD such as \texttt{com} or a country-code TLD such
as \texttt{pk}. Global analysis applies no scope filter to primary records and
uses results tagged \texttt{all}. Configured country views map each displayed
country to its ccTLD and apply that value consistently to primary queries and
materialized results. These scopes are domain-based classifications and do not
establish server location, issuer jurisdiction, or certificate-subject
geography.

\section{Indexing, Caching, and Request Processing}
\label{sec:index-cache-strategy}

Indexes support the principal access patterns for scope, domain, certificate
identity, public-key fingerprints, validity dates, cryptographic attributes,
issuer information, SANs, certificate-quality flags, and leaf
classification. Analytical result documents are indexed by scope and by the
feature-specific fields required for sorting and filtering. Frequently
requested aggregate views are served from scope-tagged, precomputed analytical
results, whereas certificate search, selected metrics, and fallback paths use
indexed primary records. Pagination and capped sample-identifier lists keep
the returned data appropriately bounded.

For applicable request paths, the optional Redis cache stores serialized
responses under keys derived from the active scope and request parameters. A
cache hit avoids repeated data access; on a miss, or when Redis is unavailable,
the service obtains the required information from the primary certificate data
or the materialized analytical results and may cache the resulting response.
Caching therefore improves repeated access without becoming a requirement for
correct operation.

A dashboard request supplies the selected scope and relevant filters, which
are applied before the service selects the appropriate data-access path. The
result is structured and returned to the dashboard, as summarized in
Figure~\ref{fig:request-sequence}.

\begin{figure}[H]
  \centering
  \begin{tikzpicture}[
    font=\scriptsize,
    >=Latex,
    participant/.style={draw, rounded corners=2pt, align=center, text width=1.95cm, minimum height=0.75cm, fill=black!5},
    life/.style={densely dashed, black!60},
    call/.style={-{Latex[length=2mm]}, line width=0.5pt},
    return/.style={-{Latex[length=2mm]}, line width=0.5pt, dashed},
    msg/.style={fill=white, inner sep=1.2pt, align=center},
    branch/.style={draw=black!45, densely dashed, rounded corners=2pt}
  ]
    \node[participant] (ui) at (0,0) {Dashboard};
    \node[participant] (mw) at (2.55,0) {Scope handling};
    \node[participant] (ctrl) at (5.10,0) {Application service};
    \node[participant] (cache) at (7.65,0) {Optional cache};
    \node[participant] (primary) at (10.20,0) {Primary certificate\\data};
    \node[participant] (results) at (12.75,0) {Materialized\\analytical results};

    \foreach \n in {ui,mw,ctrl,cache,primary,results}{\draw[life] (\n.south) -- ++(0,-7.55);}

    \draw[call] (0,-0.90) -- node[msg,above] {request, filters, and scope} (2.55,-0.90);
    \draw[call] (2.55,-1.60) -- node[msg,above] {scoped request} (5.10,-1.60);

    \draw[branch] (4.72,-1.98) rectangle (8.08,-3.28);
    \node[anchor=west, fill=white, inner sep=1pt, font=\scriptsize\itshape] at (4.86,-1.98) {cache-enabled paths};
    \draw[call] (5.10,-2.38) -- node[msg,above] {scope-aware lookup} (7.65,-2.38);
    \draw[return] (7.65,-3.00) -- node[msg,above] {hit or miss} (5.10,-3.00);

    \draw[branch] (4.72,-3.55) rectangle (13.22,-6.23);
    \node[anchor=west, fill=white, inner sep=1pt, font=\scriptsize\itshape] at (4.86,-3.55) {cache miss or uncached request; source selected as required};
    \draw[call] (5.10,-4.12) -- node[msg,above] {indexed live query} (10.20,-4.12);
    \draw[return] (10.20,-4.78) -- node[msg,above] {certificate records} (5.10,-4.78);
    \draw[call] (5.10,-5.36) -- node[msg,above] {scoped result query} (12.75,-5.36);
    \draw[return] (12.75,-5.96) -- node[msg,above] {analytical result} (5.10,-5.96);

    \draw[call] (5.10,-6.67) -- node[msg,above] {store result, if enabled} (7.65,-6.67);
    \draw[return] (5.10,-7.30) -- node[msg,above] {structured response} (0,-7.30);
  \end{tikzpicture}
  \caption{Scoped dashboard request and response sequence}
  \label{fig:request-sequence}
\end{figure}

\section{Deployment and Architectural Boundaries}
\label{sec:architectural-constraints}

The developed platform is configured as a locally operated research system.
MongoDB provides the primary persistent data store, Redis serves as an optional
response cache, and Certificate Transparency observation and certificate-parsing
processes operate as supporting components. The deployment model does not
assume a production cluster or multi-tenant operation.

Several architectural boundaries follow from this model. Scope selection is
maintained as shared application state, so strongly isolated concurrent users
would require additional request-level separation. Certificate acquisition
intentionally bypasses ordinary trust-chain and hostname verification because
its purpose is to observe certificates presented by TLS endpoints rather than
to establish browser-like trust. Dataset updates and analytical refreshes occur
in separate processing stages, with recovery handled within individual stages
rather than by one atomic transaction spanning the complete workflow. These
limitations and their implications are discussed further in
Chapter~\ref{chap:verification-security-discussion}.
